\subsection{Analysis von Sensitivity}
In this section I will analyse how robust and stable the results are, by changing the hyperparameters in the methods.
It's important, that here is just written the methodology. 
So I will concentrate on the how not the why. 
Therefor it is a part of the experimental design not the result section.

- Introduction 
    - Motivation: Why is it imporant to test the sensitivity?
    - Robustness: 
    - Stability: 
- MCD
    - Whats are the hyperparameter - support fraction (h)
    - contamination parameter ?
- k-NN  
    Hyperparameter: k 
    Hyperparametertuning: Maybe use Octuna as a pyhton framework
- InterCont
    Hyperparameter: u 


\subsection{k-NN: Choice of hyperparameter k, hyperparameter tuning}
\subsubsection{k-NN: Choice of hyperparameter k, hyperparameter tuning}
I tried to reproduce the F1-Score for all datasets except Abalone, KDDCUP99-Rev and KDDCUP99, because where not in the official ODDS Dataset. 
In the original Dataset there were sort of the same datasets, but I think the auther updated it. So I rejected them from further analysis. 
First I tried reproducing by using the classic Python machine leaning library sckit-learn. 
I noticed that the result differ from the ones in the paper.
As the auther of the paper noticed, they used the python library PyOD, wich is well known in the scientif community.
After implementing the code, I recognized that again the result differ substantially. 
Unfortunately in the reproduction package supplied by the authers, there are no description of reproducting k-NN.
They concentrated on their own .
For the F1 score of k-NN I suppose that they optimized for the best contmination.
Contamination is PyOD specific hyperparameter, that manages the sensitivity.
For every contamination rate = x, take x\% of the highest anomaliescore and mark it as one.
The difference to 100 are inliers.
Therefor I decided to use  ROCAUC wich is not influenced by contamination. 

\subsection{Breakdown-Point h of FAST-MCD}
\subsection{InterCont: hyperparameter k, m, u, hyperparameter tuning,}
In paper it is described how they did it. They call it ablation analysis.
First they switched the tanh activation function to LeakyReLU.
Then they just normalized the query and the vectors once, then they deployed a variant that just did the second normalization (the conventional way).
a variant in wich the order of the two normalizations is reversed. 
A variatnt with no normalizations. 


\subsection{Literatur}
\citep[p.5]{hilal2022financial}